% =====================================================================
% Section 7 body: Validation Against Published Computational Fluid Dynamics
% Fragment \input INSIDE \section{...}; do NOT redeclare the section.
% Numbers authoritative per _DRAFTING_BRIEF.md (brief wins over JSR_PAPER.md).
% =====================================================================

No in-house computational fluid dynamics (CFD) campaign was undertaken; this is a
deliberate scoping decision, revisited in the limitations discussion. In its place
the present method is compared against four independent published CFD studies that
together span two reference geometries, two distinct aerodynamic quantities, and
three Mach bands. The four comparators are: the unsteady Reynolds-averaged
Navier--Stokes (URANS) study of the Army--Navy Basic Finner by
Bunescu et~al.~\cite{bunescu2025}; the thin-layer Navier--Stokes computation of a
secant-ogive-cylinder-boattail projectile by Sahu, Nietubicz, and
Steger~\cite{sahu1983}; the Menter shear-stress-transport (SST) $k$--$\omega$
solution of the AGARD Model~B calibration standard by
Vidanovi\'c et~al.~\cite{vidanovic2014}; and the ANSYS Fluent pitch-damping
computations of the Basic Finner by Sznajder~\cite{sznajder2025}. A second
independent CFD source on Basic Finner pitch damping, by Bhagwandin and
Sahu~\cite{bhagwandin2013}, is used to corroborate the Sznajder finding. The
inventory is summarized in Table~\ref{tab:cfd_inventory}, and the CFD-side
panels are collected in Fig.~\ref{fig:cfd_panels}.

\begin{table*}[htbp]
  \centering
  \caption{Published-CFD comparator inventory.}
  \label{tab:cfd_inventory}
  \begin{tabularx}{\textwidth}{@{}l l l l X@{}}
    \toprule
    Source & Geometry & Quantity & Mach range & Comparison status \\
    \midrule
    Bunescu et~al.~\cite{bunescu2025} (URANS)
      & Basic Finner
      & $C_X$
      & 0.4--3.5
      & Comparator wired; trend correct, loose absolute agreement \\
    Sahu et~al.~\cite{sahu1983} (thin-layer NS)
      & Ogive-cylinder-boattail
      & $C_{D,b}$, $C_{D}$
      & 0.9--1.2
      & Memo only; digitization deferred to future work \\
    Vidanovi\'c et~al.~\cite{vidanovic2014} (SST $k$--$\omega$)
      & AGARD Model~B
      & $C_D$, $C_L$, $C_m$
      & 0.596, 1.602
      & Reference dataset; no AGARD-B build shipped \\
    Sznajder~\cite{sznajder2025} (Fluent)
      & Basic Finner
      & $C_{m_q}+C_{m_{\dot\alpha}}$
      & 0.9--5.0
      & Comparator wired; supersonic MAPE 31.6\% ($n=8$, $M\ge1.29$) \\
    Bhagwandin \& Sahu~\cite{bhagwandin2013} (Fluent)
      & Basic Finner
      & $C_{m_q}+C_{m_{\dot\alpha}}$
      & 0.6--4.5
      & Second-source corroboration; MAPE 50.8\% ($n=13$) \\
    \bottomrule
  \end{tabularx}
\end{table*}

\subsection{Basic Finner Static Axial Force versus Bunescu URANS}

Bunescu et~al.~\cite{bunescu2025} reported URANS predictions of normal- and
axial-force coefficients on the standard four-rectangular-fin Army--Navy Basic
Finner at $M = 0.4$, 0.95, 1.6, 2.5, and 3.5, using a realizable $k$--$\varepsilon$
closure with corroborating SST $k$--$\omega$ comparisons. The present method was
exercised on the identical geometry at the same five Mach numbers and at zero angle
of attack. Absolute agreement in the axial-force coefficient $C_X$ is loose: the
largest residuals occur near the transonic point $M = 0.95$, where the
shock-induced separation resolved by the URANS field is not represented by the
engineering pressure-drag model, and at the supersonic points, where the
semi-empirical viscous treatment differs from the Reynolds-number-resolved
boundary layer. The Mach trend, however, is reproduced correctly: a transonic rise
to a peak near $M \approx 1$ followed by monotone decay through $M = 3.5$. This
comparison is reported as a qualitative trend check rather than as a quantitative
benchmark.

\subsection{Ogive-Cylinder-Boattail Base Drag versus Sahu Thin-Layer Navier--Stokes}

Sahu, Nietubicz, and Steger~\cite{sahu1983} computed base, pressure, friction, and
total drag for a secant-ogive-cylinder-boattail projectile across the transonic
band $M = 0.9$--1.2 at $\mathrm{Re}_L = 4.5\times10^{6}$, overlaying their
thin-layer Navier--Stokes solution against sting-mounted wind-tunnel data and a
semi-empirical correlation. This triple comparison would provide an ideal anchor
for a transonic base-drag panel against the present method's base-drag and
pressure-drag pipelines. The digitization of this dataset was deferred during
preparation of the present manuscript; the comparison is therefore retained as
near-term future work rather than presented as a result, and the corresponding
panel in Fig.~\ref{fig:cfd_panels} is flagged accordingly.

\subsection{AGARD Model~B Reference versus Vidanovi\'c SST $k$--$\omega$}

Vidanovi\'c et~al.~\cite{vidanovic2014} reported ANSYS Fluent SST $k$--$\omega$
predictions of total drag, lift, and pitching-moment coefficients on the AGARD
Model~B calibration standard at $M = 0.596$ and $M = 1.602$ over an
angle-of-attack sweep, validated against wind-tunnel data from the VTI~T-38
trisonic facility. The published CFD-versus-experiment agreement is within a few
percent in $C_D$ and below one percent in $C_L$ over the test envelope. The present
method does not ship an AGARD-B build---the equilateral-triangle delta wing with a
biconvex section lies at the edge of the validity of the underlying fin-set
model---so no quantitative comparison is reported here. Consistent with the
benchmark roll-up, the AGARD-B configuration is treated as a qualitative trend
reference only (the corresponding total-drag agreement is loose, with mean absolute
percentage error of order $22\%$), and the Vidanovi\'c CFD is displayed in
Fig.~\ref{fig:cfd_panels} as a reference dataset against which a future comparator
build can be benchmarked. The omission is acknowledged as a deferred future-work
item.

\subsection{Basic Finner Pitch Damping versus Sznajder Fluent, Corroborated by Bhagwandin and Sahu}
\label{sec:cfd_cmq}

The most informative CFD comparison concerns the pitch-damping derivative on the
Basic Finner. Sznajder~\cite{sznajder2025} computed $C_{m_q}$ and
$C_{m_{\dot\alpha}}$ using three independent CFD techniques---a steady
moving-reference-frame method, a dynamic-mesh forced-oscillation method, and a
step-perturbation indicial-response method---over $M = 0.9$--5.0. The three methods
agreed to within approximately $3\%$ of one another and were independently
validated against free-flight experimental data. The present method exposes the
experimentally observable damping sum $C_{m_q}+C_{m_{\dot\alpha}}$, which was
compared against the Sznajder solution on a common grid with the following two
findings.

\begin{itemize}
  \item \emph{Supersonic band}, $M = 1.29$--4.5 ($n = 8$ points): the present
        method underpredicts the magnitude of the damping sum by $27$ to $36\%$,
        with the sign and Mach trend correct. The mean absolute percentage error
        over this band is $31.6\%$.
  \item \emph{Transonic peak}, $M \approx 1.08$--1.11 ($n = 2$ points): the present
        method overshoots the magnitude of the damping sum by $110$ to $160\%$;
        the CFD does not exhibit a comparable peak in the sum.
\end{itemize}

The transonic overshoot is traced to the Gaussian augmentation applied to the
strip-theory damping near $M = 1$ in the stability model; the supersonic
underprediction reflects a constant-factor bias in the strip-theory damping
coefficient. Bhagwandin and Sahu~\cite{bhagwandin2013} provide an independent ANSYS
Fluent reference for the same geometry over $M = 0.6$--4.5, validated against the
same free-flight dataset; their comparator yields a mean absolute percentage error
of $50.8\%$ over thirteen points, with the same sign and the same direction of the
residual.

Two independent CFD sources therefore converge on the same two findings: a
$27$--$36\%$ supersonic underprediction of pitch damping (with correct sign and
trend) and a transonic-peak over-augmentation. Consistent with this evidence, the
transonic pitch-damping augmentation is held to be a B-level model element rather
than a validated one. Both findings are taken up explicitly in the limitations
discussion.

\begin{figure*}[htbp]
  \centering
  \includegraphics[width=\textwidth]{cfd_validation_panels}
  \caption{Published-CFD comparator composite. Panel~A: Basic Finner axial-force
  coefficient versus the Bunescu URANS solution~\cite{bunescu2025} (correct trend,
  loose absolute agreement). Panel~B: ogive-cylinder-boattail base drag versus the
  Sahu thin-layer Navier--Stokes solution~\cite{sahu1983} (digitization deferred).
  Panel~C: AGARD Model~B reference dataset versus the Vidanovi\'c SST $k$--$\omega$
  solution and VTI~T-38 experiment~\cite{vidanovic2014}. Panel~D: Basic Finner
  pitch-damping sum versus the Sznajder Fluent solution~\cite{sznajder2025},
  corroborated by Bhagwandin and Sahu~\cite{bhagwandin2013}.}
  \label{fig:cfd_panels}
\end{figure*}
